\documentclass[10pt]{article}
\usepackage[T1]{fontenc}
\usepackage[utf8]{inputenc}
\usepackage[french]{babel}
\usepackage{amsmath,amssymb,mathrsfs,mathpazo}
\usepackage{geometry}
\geometry{a4paper,margin=2cm}
\usepackage{graphicx,array,tabularx,tabularray,colortbl,multirow,multicol,booktabs,enumitem,pifont,ulem,fancybox,diagbox,cancel}
\usepackage{pstricks,pstricks-add,pst-plot,pst-tree,pst-node,pst-text,pst-func,pst-math,pst-bspline,pst-3dplot,pst-eucl,pst-solides3d}
\let\pscalculate\undefined
\usepackage{pst-bezier}
\usepackage{tikz,pgfplots,tkz-tab}
\usetikzlibrary{arrows,arrows.meta,patterns,calc,positioning,trees,decorations.markings,shapes.geometric,shapes.misc,shapes.arrows,intersections,decorations.pathmorphing,decorations.pathreplacing,backgrounds,shadows,babel}
\pgfplotsset{compat=1.18}
\usepgfplotslibrary{fillbetween}
\usepackage[np]{numprint}
\usepackage{listings,makecell,pas-tableur,twoopt}
\let\barre\undefined
\usepackage[pstricks]{tablvar}
\definecolor{darkgreen}{RGB}{0,100,0}
\providecommand{\red}{\color{red}}
\providecommand{\blue}{\color{blue}}
\providecommand{\green}{\color{green}}
\providecommand{\black}{\color{black}}
\providecommand{\white}{\color{white}}
\providecommand{\gray}{\color{gray}}
\providecommand{\cyan}{\color{cyan}}
\providecommand{\magenta}{\color{magenta}}
\providecommand{\yellow}{\color{yellow}}
\providecommand{\violet}{\color{violet}}
\providecommand{\orange}{\color{orange}}
\providecommand{\euro}{\texteuro}
\providecommand{\EUR}{\texteuro}
\providecommand{\ieme}{\textsuperscript{e}}
\providecommand{\degres}{\ensuremath{^\circ}}
\providecommand{\degree}{\ensuremath{^\circ}}
\providecommand{\vv}[1]{\overrightarrow{#1}}
\providecommand{\Coord}[1]{#1}
\providecommand{\decosix}{}
\providecommand{\starredbullet}{\ensuremath{\bullet}}
\pagestyle{empty}
\setlength{\parindent}{0pt}
\usepackage{adjustbox}\DeclareUnicodeCharacter{2236}{:}\DeclareUnicodeCharacter{2190}{\ensuremath{\leftarrow}}
\providecommand{\R}{\mathbb{R}}
\providecommand{\N}{\mathbb{N}}
\providecommand{\Z}{\mathbb{Z}}
\providecommand{\Q}{\mathbb{Q}}
\providecommand{\D}{\mathbb{D}}
\providecommand{\C}{\mathbb{C}}
\providecommand{\e}{{\mathrm{e}}}
\providecommand{\vect}[1]{\overrightarrow{#1}}
\providecommand{\vectt}[1]{\overrightarrow{\mathrm{#1}}}
\providecommand{\vvt}[1]{\overrightarrow{\mathrm{#1}}}
\providecommand{\np}[2][]{{#2~#1}}
\providecommand{\vv}[1]{{\overrightarrow{#1}}}
\renewcommand{\d}{{\mathrm{d}}}
\providecommand{\degres}{{^\circ}}
\providecommand{\degre}{{^\circ}}
\providecommand{\blue}{\color{blue}}
\providecommand{\red}{\color{red}}
\providecommand{\enskip}{\quad}
\providecommand{\llbracket}{\lbrack\!\lbrack}
\providecommand{\rrbracket}{\rbrack\!\rbrack}
\providecommand{\euro}{{\text{€}}}
\providecommand{\no}{{\text{n°}}}
\providecommand{\interval}[2]{{[#1;#2]}}
\providecommand{\intervFO}[2]{{[#1;#2[}}
\providecommand{\dotfill}{\ldots}

\providecommand{\ds}{\displaystyle}
\providecommand{\textding}[1]{#1}
\providecommand{\Oij}{\(\mathrm{O};\vec{\imath},\vec{\jmath}\)}
\providecommand{\Oijk}{\(\mathrm{O};\vec{\imath},\vec{\jmath},\vec{k}\)}
\providecommand{\Ouv}{\(\mathrm{O};\vec{u},\vec{v}\)}
\providecommand{\intervFF}[2]{\ensuremath{[#1;#2]}}
\renewcommand{\euro}{\ifmmode\text{\texteuro}\else\texteuro\fi}
\DeclareUnicodeCharacter{00B0}{\ifmmode^{\circ}\else\textdegree\fi}
\DeclareUnicodeCharacter{20AC}{\ifmmode\text{\texteuro}\else\texteuro\fi}
% Énoncés : Éducation nationale. Transcription et corrigés : APMEP, Denis Vergès et contributeurs. Découpage et classement : CoopMaths, Sébastien Lozano, Éric Elter et contributeurs. Adaptation de la présentation : Maths Explorer.
% CC BY-SA : https://coopmaths.fr/www/about/
% https://coopmaths.fr/alea/static/bac/2023/tex/bac_2023_09_sujet1_metropole_4.tex
\begin{document}
\noindent\textit{APMEP / CoopMaths. CC BY-SA. Présentation adaptée par Maths Explorer.}\par\medskip
\medskip

\textbf{PARTIE A}

\medskip

On définit sur l'intervalle $]0~;~+\infty[$ la fonction $g$ par : 
\[g(x) = \dfrac 2x - \dfrac{1}{x^2} + \ln x \quad \text{où ln désigne la fonction logarithme népérien}.\]

On admet que la fonction $g$ est dérivable sur $]0~;~+\infty[ = I$ et on note $g'$ sa fonction dérivée.

\medskip

\begin{enumerate}
\item Montrer que pour $x > 0$, le signe de $g'(x)$ est celui du trinôme du second degré $\left(x^2 - 2x + 2\right)$.
\item En déduire que la fonction $g$ est strictement croissante sur $]0~;~+\infty[$.
\item Montrer que l'équation $g(x) = 0$ admet une unique solution sur l'intervalle [0,5~;~1], que l'on notera $\alpha$.
\item On donne le tableau de signes de $g$ sur l'intervalle $]0~;~+\infty[ = I$ :
\begin{center}
\psset{unit=1cm}
\begin{pspicture}(4,1)
\psline(0,0.5)(4,0.5)\psline(1,0)(1,1)\psline(1.1,0)(1.1,0.5)
\uput[u](0.5,0.4){$x$} \uput[u](1.2,0.4){$0$} \uput[u](2.5,0.4){$\alpha$} \uput[u](3.7,0.4){$+\infty$}
\uput[u](0.5,-0.1){$g(x)$}\uput[u](1.75,-0.1){$-$}\uput[u](2.5,-0.1){$0$}
\uput[u](3.25,-0.1){$+$}
\end{pspicture}
\end{center}

Justifier ce tableau de signes à l'aide des résultats obtenus aux questions précédentes.
\end{enumerate}

\medskip

\textbf{PARTIE B}

\medskip

On considère la fonction $f$ définie sur l'intervalle $]0~;~+\infty[ = I$ par: 
\[f(x) = \e^x \ln x.\]

On note $\mathcal{C}_f$ la courbe représentative de $f$ dans un repère orthonormé.

\medskip

\begin{enumerate}
\item On admet que la fonction $f$ est deux fois dérivable sur $]0~;~+\infty[$, on note $f'$ sa fonction dérivée, $f''$ sa fonction dérivée seconde et on admet que :

pour tout nombre réel $x > 0,\:\: f'(x) = \e^x \left(\dfrac 1x + \ln x\right)$.

Démontrer que, pour tout nombre réel $x > 0$, on a : $f''(x) = \e^x \left(\dfrac 2x - \dfrac{1}{x^2} + \ln x\right)$.
\item On pourra remarquer que pour tout réel $x > 0$,\: $f''(x) = \e^x \times  g(x)$, où $g$ désigne la fonction étudiée dans la partie A.
\item 
	\begin{enumerate}
		\item Dresser le tableau de signes de la fonction $f''$ sur $]0~;~+\infty[$. Justifier.
		\item Justifier que la courbe $\mathcal{C}_f$ admet un unique point d'inflexion A.
		\item Étudier la convexité de la fonction $f$ sur l'intervalle $]0~;~+\infty[$. Justifier.
	\end{enumerate}
\item
	\begin{enumerate}
		\item Calculer les limites de $f$ aux bornes de son ensemble de définition.
		\item Montrer que $f'(\alpha) = \dfrac{\e^{\alpha}}{\alpha^2}(1 - \alpha)$.

On rappelle que $\alpha$ est l'unique solution de l'équation $g(x) = 0$.
		\item Démontrer que $f'(\alpha) > 0$ et en déduire le signe de $f'(x)$ pour $x$ appartenant à $]0~;~+\infty[$.
		\item En déduire le tableau de variations complet de la fonction $f$ sur $]0~;~+\infty[$.
	\end{enumerate}
\end{enumerate}

\end{document}
